\section{Full Repositories List}

\resizebox{\textwidth}{!}{
\begin{tabular}{llrrrr}
\toprule
Type & Repository & License & \#Instances & \#Files & LoC \\
\midrule
\multirow{26}{*}{AI/ML} & \href{https://github.com/deepset-ai/haystack}{deepset-ai/haystack} & Apache-2.0 & 64 & 433 & 84.8k \\
& \href{https://github.com/instructlab/instructlab}{instructlab/instructlab} & Apache-2.0 & 52 & 142 & 28.2k \\
& \href{https://github.com/keras-team/keras}{keras-team/keras} & Apache-2.0 & 48 & 900 & 249.7k \\
& \href{https://github.com/kedro-org/kedro}{kedro-org/kedro} & Apache-2.0 & 27 & 179 & 40.4k \\
& \href{https://github.com/pytorch/torchtune}{pytorch/torchtune} & BSD-3-Clause & 14 & 448 & 92.4k \\
& \href{https://github.com/jupyterlab/jupyter-ai}{jupyterlab/jupyter-ai} & BSD-3-Clause & 13 & 81 & 9.0k \\
& \href{https://github.com/run-llama/llama\_deploy}{run-llama/llama\_deploy} & MIT & 12 & 216 & 15.1k \\
& \href{https://github.com/stanfordnlp/dspy}{stanfordnlp/dspy} & MIT & 10 & 222 & 30.6k \\
& \href{https://github.com/projectmesa/mesa}{projectmesa/mesa} & Apache-2.0 & 9 & 109 & 20.3k \\
& \href{https://github.com/huggingface/smolagents}{huggingface/smolagents} & Apache-2.0 & 5 & 65 & 21.4k \\
& \href{https://github.com/theOehrly/Fast-F1}{theOehrly/Fast-F1} & MIT & 4 & 92 & 20.7k \\
& \href{https://github.com/cyclotruc/gitingest}{cyclotruc/gitingest} & MIT & 3 & 39 & 4.7k \\
& \href{https://github.com/modelcontextprotocol/python-sdk}{modelcontextprotocol/python-sdk} & MIT & 2 & 114 & 13.4k \\
& \href{https://github.com/camel-ai/camel}{camel-ai/camel} & Apache-2.0 & 2 & 799 & 130.1k \\
& \href{https://github.com/hiyouga/LLaMA-Factory}{hiyouga/LLaMA-Factory} & Apache-2.0 & 2 & 170 & 31.2k \\
& \href{https://github.com/feast-dev/feast}{feast-dev/feast} & Apache-2.0 & 2 & 673 & 103.3k \\
& \href{https://github.com/openai/openai-agents-python}{openai/openai-agents-python} & MIT & 1 & 212 & 29.8k \\
& \href{https://github.com/huggingface/datasets}{huggingface/datasets} & Apache-2.0 & 1 & 207 & 69.6k \\
& \href{https://github.com/stanford-crfm/helm}{stanford-crfm/helm} & Apache-2.0 & 1 & 891 & 122.1k \\
& \href{https://github.com/freqtrade/freqtrade}{freqtrade/freqtrade} & GPL-3.0 & 1 & 458 & 130.1k \\
& \href{https://github.com/lss233/kirara-ai}{lss233/kirara-ai} & AGPL-3.0 & 1 & 261 & 25.5k \\
& \href{https://github.com/arviz-devs/arviz}{arviz-devs/arviz} & Apache-2.0 & 1 & 259 & 50.6k \\
& \href{https://github.com/qubvel-org/segmentation\_models.pytorch}{qubvel-org/segmentation\_models.pytorch} & MIT & 1 & 130 & 18.6k \\
& \href{https://github.com/scikit-learn-contrib/category\_encoders}{scikit-learn-contrib/category\_encoders} & BSD-3-Clause & 1 & 71 & 12.9k \\
& \href{https://github.com/huggingface/open-r1}{huggingface/open-r1} & Apache-2.0 & 1 & 29 & 4.0k \\
& \href{https://github.com/gptme/gptme}{gptme/gptme} & MIT & 1 & 124 & 23.3k \\
\midrule
\multirow{23}{*}{DevOps} & \href{https://github.com/conan-io/conan}{conan-io/conan} & MIT & 136 & 1056 & 162.5k \\
& \href{https://github.com/pylint-dev/pylint}{pylint-dev/pylint} & GPL-2.0 & 57 & 2301 & 116.8k \\
& \href{https://github.com/sphinx-doc/sphinx}{sphinx-doc/sphinx} & N/A & 39 & 718 & 140.3k \\
& \href{https://github.com/pdm-project/pdm}{pdm-project/pdm} & MIT & 34 & 221 & 32.3k \\
& \href{https://github.com/beeware/briefcase}{beeware/briefcase} & BSD-3-Clause & 24 & 508 & 89.3k \\
& \href{https://github.com/bridgecrewio/checkov}{bridgecrewio/checkov} & Apache-2.0 & 21 & 4551 & 234.9k \\
& \href{https://github.com/joke2k/faker}{joke2k/faker} & MIT & 20 & 754 & 351.4k \\
& \href{https://github.com/python-attrs/attrs}{python-attrs/attrs} & MIT & 10 & 52 & 18.6k \\
& \href{https://github.com/ipython/ipython}{ipython/ipython} & BSD-3-Clause & 10 & 293 & 79.8k \\
& \href{https://github.com/koxudaxi/datamodel-code-generator}{koxudaxi/datamodel-code-generator} & MIT & 10 & 599 & 60.3k \\
& \href{https://github.com/tox-dev/tox}{tox-dev/tox} & MIT & 7 & 225 & 23.8k \\
& \href{https://github.com/dynaconf/dynaconf}{dynaconf/dynaconf} & MIT & 6 & 463 & 55.1k \\
& \href{https://github.com/pypa/twine}{pypa/twine} & Apache-2.0 & 6 & 34 & 6.6k \\
& \href{https://github.com/wemake-services/wemake-python-styleguide}{wemake-services/wemake-python-styleguide} & MIT & 6 & 396 & 52.5k \\
& \href{https://github.com/Delgan/loguru}{Delgan/loguru} & MIT & 6 & 168 & 19.2k \\
& \href{https://github.com/kubernetes-client/python}{kubernetes-client/python} & Apache-2.0 & 3 & 783 & 267.2k \\
& \href{https://github.com/olofk/fusesoc}{olofk/fusesoc} & BSD-2-Clause & 2 & 45 & 8.8k \\
& \href{https://github.com/amoffat/sh}{amoffat/sh} & MIT & 2 & 5 & 7.4k \\
& \href{https://github.com/facebookresearch/hydra}{facebookresearch/hydra} & MIT & 2 & 439 & 41.4k \\
& \href{https://github.com/home-assistant/supervisor}{home-assistant/supervisor} & Apache-2.0 & 2 & 541 & 82.3k \\
& \href{https://github.com/FreeOpcUa/opcua-asyncio}{FreeOpcUa/opcua-asyncio} & LGPL-3.0 & 1 & 168 & 344.4k \\
& \href{https://github.com/pytest-dev/pytest}{pytest-dev/pytest} & MIT & 1 & 260 & 99.5k \\
& \href{https://github.com/iterative/dvc}{iterative/dvc} & Apache-2.0 & 1 & 554 & 85.3k \\
\midrule
\multirow{12}{*}{Web} & \href{https://github.com/reflex-dev/reflex}{reflex-dev/reflex} & Apache-2.0 & 44 & 376 & 89.8k \\
& \href{https://github.com/sissbruecker/linkding}{sissbruecker/linkding} & MIT & 29 & 193 & 26.4k \\
& \href{https://github.com/Kozea/WeasyPrint}{Kozea/WeasyPrint} & BSD-3-Clause & 19 & 144 & 70.0k \\
& \href{https://github.com/python-telegram-bot/python-telegram-bot}{python-telegram-bot/python-telegram-bot} & GPL-3.0 & 16 & 464 & 140.8k \\
& \href{https://github.com/python-babel/babel}{python-babel/babel} & BSD-3-Clause & 11 & 75 & 23.1k \\
& \href{https://github.com/falconry/falcon}{falconry/falcon} & Apache-2.0 & 11 & 262 & 58.5k \\
& \href{https://github.com/aiogram/aiogram}{aiogram/aiogram} & MIT & 11 & 861 & 69.8k \\
& \href{https://github.com/privacyidea/privacyidea}{privacyidea/privacyidea} & AGPL-3.0 & 10 & 483 & 167.5k \\
& \href{https://github.com/urllib3/urllib3}{urllib3/urllib3} & MIT & 10 & 81 & 31.3k \\
& \href{https://github.com/ag2ai/faststream}{ag2ai/faststream} & Apache-2.0 & 6 & 1267 & 85.1k \\
& \href{https://github.com/encode/starlette}{encode/starlette} & BSD-3-Clause & 5 & 66 & 17.2k \\
& \href{https://github.com/scrapinghub/dateparser}{scrapinghub/dateparser} & BSD-3-Clause & 2 & 274 & 67.1k \\
\bottomrule
\end{tabular}
}


%%%

\resizebox{\textwidth}{!}{
\begin{tabular}{llrrrr}
\toprule
Type & Repository & License & \#Instances & \#Files & LoC \\
\midrule
\multirow{5}{*}{Web} & \href{https://github.com/pallets/flask}{pallets/flask} & BSD-3-Clause & 2 & 83 & 17.8k \\
& \href{https://github.com/scrapy-plugins/scrapy-splash}{scrapy-plugins/scrapy-splash} & BSD-3-Clause & 1 & 25 & 3.4k \\
& \href{https://github.com/psf/requests}{psf/requests} & Apache-2.0 & 1 & 36 & 11.2k \\
& \href{https://github.com/jpadilla/pyjwt}{jpadilla/pyjwt} & MIT & 1 & 26 & 6.9k \\
& \href{https://github.com/slackapi/bolt-python}{slackapi/bolt-python} & MIT & 1 & 562 & 60.8k \\
\midrule
\multirow{8}{*}{Database} & \href{https://github.com/pydata/xarray}{pydata/xarray} & Apache-2.0 & 29 & 226 & 179.2k \\
& \href{https://github.com/geopandas/geopandas}{geopandas/geopandas} & BSD-3-Clause & 21 & 87 & 47.3k \\
& \href{https://github.com/reata/sqllineage}{reata/sqllineage} & MIT & 18 & 103 & 9.7k \\
& \href{https://github.com/patroni/patroni}{patroni/patroni} & MIT & 17 & 117 & 45.9k \\
& \href{https://github.com/piskvorky/smart\_open}{piskvorky/smart\_open} & MIT & 6 & 64 & 12.4k \\
& \href{https://github.com/wireservice/csvkit}{wireservice/csvkit} & MIT & 3 & 48 & 6.6k \\
& \href{https://github.com/jazzband/tablib}{jazzband/tablib} & MIT & 2 & 32 & 6.6k \\
& \href{https://github.com/Flexget/Flexget}{Flexget/Flexget} & MIT & 2 & 657 & 108.6k \\
\midrule
\multirow{8}{*}{Scientific} & \href{https://github.com/pvlib/pvlib-python}{pvlib/pvlib-python} & BSD-3-Clause & 29 & 178 & 59.9k \\
& \href{https://github.com/python-control/python-control}{python-control/python-control} & BSD-3-Clause & 15 & 155 & 70.7k \\
& \href{https://github.com/mikedh/trimesh}{mikedh/trimesh} & MIT & 14 & 248 & 74.8k \\
& \href{https://github.com/PyPSA/PyPSA}{PyPSA/PyPSA} & MIT & 10 & 129 & 32.3k \\
& \href{https://github.com/shapely/shapely}{shapely/shapely} & BSD-3-Clause & 9 & 158 & 34.0k \\
& \href{https://github.com/pybamm-team/PyBaMM}{pybamm-team/PyBaMM} & BSD-3-Clause & 6 & 581 & 113.4k \\
& \href{https://github.com/beancount/beancount}{beancount/beancount} & GPL-2.0 & 2 & 194 & 48.0k \\
& \href{https://github.com/sympy/sympy}{sympy/sympy} & N/A & 2 & 1574 & 760.5k \\
\midrule
\multirow{4}{*}{CLI} & \href{https://github.com/streamlink/streamlink}{streamlink/streamlink} & BSD-2-Clause & 39 & 510 & 84.0k \\
& \href{https://github.com/beetbox/beets}{beetbox/beets} & MIT & 9 & 193 & 69.3k \\
& \href{https://github.com/yt-dlp/yt-dlp}{yt-dlp/yt-dlp} & Unlicense & 5 & 1177 & 244.6k \\
& \href{https://github.com/jarun/buku}{jarun/buku} & GPL-3.0 & 2 & 27 & 7.1k \\
\midrule
\multirow{3}{*}{Misc} & \href{https://github.com/matplotlib/matplotlib}{matplotlib/matplotlib} & N/A & 85 & 904 & 263.6k \\
& \href{https://github.com/fonttools/fonttools}{fonttools/fonttools} & MIT & 12 & 512 & 192.6k \\
& \href{https://github.com/pytransitions/transitions}{pytransitions/transitions} & MIT & 2 & 39 & 12.4k \\
\midrule
\multirow{2}{*}{Cloud} & \href{https://github.com/aws-cloudformation/cfn-lint}{aws-cloudformation/cfn-lint} & MIT-0 & 102 & 2422 & 160.2k \\
& \href{https://github.com/icloud-photos-downloader/icloud\_photos\_downloader}{icloud-photos-downloader/icloud\_photos\_downloader} & MIT & 4 & 73 & 15.5k \\
\midrule
\multirow{2}{*}{Desktop} & \href{https://github.com/qtile/qtile}{qtile/qtile} & MIT & 6 & 405 & 81.6k \\
& \href{https://github.com/pwr-Solaar/Solaar}{pwr-Solaar/Solaar} & GPL-2.0 & 3 & 94 & 33.7k \\
\bottomrule
\end{tabular}
}

\section{Task Formulation}
\label{sec:preliminary}

\begin{figure}[!h]
    \centering
    \includegraphics[width=0.8\linewidth]{figures/task.pdf}
    \caption{The issue-resolving task requires the model to generate a patch that addresses a given issue, with its correctness evaluated through test execution.}
    \label{fig:task}
\end{figure}

Issue resolving is the task that introduced by SWE-bench~\cite{swebench} for benchmarking AI coding capabilities. In simple terms, it simulates the process of a developer submitting a pull request to address an issue. The formulation of the issue-resolving task is illustrated in Figure~\ref{fig:task}.

\paragraph{Generating Patch.} The task input includes the problem statement of the issue, which is the description written by the issue reporter, as well as a snapshot of the codebase at the time the issue was filed (obtained by resetting to the \texttt{base\_commit}). The model has access to full content of the codebase, after then it is tasked with generating a patch that fixes the given issue, analogous to the file changes submitted in a pull request. In practice, the expected output is in the \texttt{.diff} format.

\paragraph{Evaluating Patch.} Once a patch is proposed by the model, we assess its correctness by applying it to the target codebase and executing the repository’s test suite. The output of the test execution are parsed using a log parser function, which extracts the status of each individual test case. These results are then compared against the expected test case transitions pre-defined for the issue, specifically \texttt{FAIL\_TO\_PASS} and \texttt{PASS\_TO\_PASS}. \texttt{FAIL\_TO\_PASS} refers to test cases that originally failed prior to patch application—typically those introduced in the corresponding pull request—and are expected to pass if the proposed solution is correct. A correct patch should successfully cause these failing tests to pass, without causing regressions in the already passing tests.

\section{Performance vs Repository difficulty}
\label{appendix:perform_vs_repo}
The following Figure~\ref{fig:perf_with_repo} plots resolved rate for every repository against its size (Python files on the x-axis, LOC on the y-axis). For detailed interpreatation of the figure, please see Section \ref{sec:perform_vs_diff}.
\begin{figure}[h!]
    \centering
    \includegraphics[width=\linewidth]{figures/repo_resolved_scatter.pdf}
    \caption{Resolved rate in relation to the number of files and lines of code of a repository. }
    \label{fig:perf_with_repo}
\end{figure}



\section{Dataset Fields}

Table~\ref{tab:fields} provides a detailed description of the fields included in the SWE-bench-Live dataset, along with how they are obtained during the curation process.

\begin{table}[h]
    \centering
    \caption{The required fields for a typical issue-solving task instance. Fields marked with * are \textbf{newly added} in SWE-bench-Live compared to SWE-bench.}
    \resizebox{\textwidth}{!}{
    \begin{tabular}{c|c|p{10cm}}
    \toprule
      \texttt{Field} & \texttt{Type} & Description\\
    \midrule
      \texttt{base\_commit} & \texttt{str} & The commit on which the pull request is based, representing the repository state before the issue is resolved. \\
      \texttt{patch} & \texttt{str} & Gold patch proposed by the pull request, in \texttt{.diff} format. \\
      \texttt{test\_patch} & \texttt{str} & Modifications to the test suite proposed by the pull request that are typically used to check whether the issue has been resolved. \\
      \texttt{problem\_statement} & \texttt{str} & Issue description text, typically describing the bug or requested feature, used as the task problem statement.\\
      \texttt{FAIL\_TO\_PASS} & \texttt{List[str]} & Test cases that are expected to successfully transition from failing to passing are used to evaluate the correctness of the patch.\\
      \texttt{PASS\_TO\_PASS} & \texttt{List[str]} & Test cases that are already passing prior to applying the gold patch. A correct patch shouldn't introduce regression failures in these tests.\\
      \texttt{*image\_key} & \texttt{str} & Instance-level docker image that provides an execution environment. \\
      \texttt{*test\_cmds} & \texttt{List[str]} & The command(s) used to run the test suite is identified by the verify agent in \method{}. It is required to enable detailed logging of each test item's status (e.g., by using the \texttt{pytest -rA} option).  \\
      \texttt{*log\_parser} & \texttt{str} & The type of log parser required for the instance—by default, \texttt{pytest}.\\
    \bottomrule
    \end{tabular}
    }
    \label{tab:fields}
\end{table}

\section{Experimental Setup Details}
\label{sec:experimental_setup}

In this section, we present additional details of the experimental setup to facilitate reproducibility.

\paragraph{Hyperparameters used in the experiments.} For OpenHands, we set a maximum of 60 iterations per instance, with the LLM configured to use a temperature of 0.0 and a top-p value of 1.0 as default. For SWE-agent, we limit the number of LLM calls to 100 per instance, with the temperature set to 0.0 and top-p to 1.0. For Agentless, both the number of localization samples and repair samples are set to 1, corresponding to a single rollout. The LLM temperature is set to 0.8 during the localization phase, as defined by the agent’s default, and 0.0 for all other phases. In our experiments, we omit the regression test-based reranking stage of Agentless, retaining only the localization and repair stages. The LLM calls within \method{} are configured with a temperature of 0.0.

\paragraph{Random seed in subset splitting.} The only stochastic component in this work arises during the sampling of the lite subset, where we set the random seed to 42.

\paragraph{Computational resources.} All LLM calls in this work are made through official APIs. The experiments involve parallel execution of multiple Docker containers for test execution. We conduct all the experiments on a CPU server equipped with an Intel Xeon Gold 6338 @ 2.00GHz (128 cores) and 2TB of RAM.

\section{Limitations}
\label{sec:limitations}
% List the limitations here
\textbf{Randomness caused by LLMs}: We use the LLMs as the core engine to conduct all the experiments, which might lead to potential randomness caused by different LLM calls. Since the experiments require extensive LLM calls while the overall budget is limited, we do not repeat the experiments for multiple times. To reduce the randomness, we control the execution environment to be the same and set the temperature and top\_p to zero. 

\textbf{Language limitation}: Our benchmark \benchmark primarily focuses on the Python language only, which might be limited. Since our key contribution is to propose a live benchmark with an automated and scalable method, we follow the same language choice as  existing benchmarks like SWE-bench. In the future, we plan to extend \benchmark to multiple languages such as Java, Go, and etc.
% \# extend to other languages

\section{Prompts in \method{}}



\label{appdix:launchprompt}

\begin{tcolorbox}[colback=gray!5,colframe=black!50,fontupper=\ttfamily,title=Prompt for Relevant Files Identification]
Given this repository structure:

------ BEGIN REPOSITORY STRUCTURE ------

\{structure\}

------ END REPOSITORY STRUCTURE ------


List the most relevant files for setting up a development environment, including:

0. CI/CD configuration files

1. README files

2. Documentation

3. Installation guides

4. Development setup guides

Format each file with its relative path (relative to project root) to be wrapped with tag <file> </file>, one per line.
\end{tcolorbox}

\begin{tcolorbox}[colback=gray!5,colframe=black!50,fontupper=\ttfamily,title=Prompt for Setup Agent]
You are a developer. Your task is to install dependencies and set up a environment that is able to run the tests of the project.\\

- You start with an initial Docker container based on \{base\_image\}.

- You interact with a Bash session inside this container.

- Project files are located in /testbed within the container, and your current working directory of bash is already set to /testbed.

- No need to clone the project again.\\

The final objective is to successfully run the tests of the project.
\\
\#\#\# Attention:

- For Python project, you should make sure the package is installed from source in the editable mode before running tests (for example 'pip install -e .') to have a development environment.

- For Python project, avoid use tox to run test if possible as it is designed specifically for CI. Read tox.ini file to find how to setup and run the test.

You run in a loop of Thought, Action, Observation.
At the end of the loop you should use Action to stop the loop.

Use Thought to describe your thoughts about the question you have been asked.

Use Action to run one of the actions available to you.

Observation will be the result of running those actions.

> Important Note: Each step, reply with only **one** (Thought, Action) pair.

> Important Note: Do not reply **Observation**, it will be provided by the system.


Your available actions are:

\{tools\}

Observation will be the result of running those actions.\\

Project Structure: the structure of the project, including files and directories.

Related Files: the content of related files of the project that may help you understand the project.

Thought: you should always think about what to do

Action: decide an action to take

Observation: the result of the action
\\

... (this Thought/Action/Observation can repeat N times) ...
\\

Thought: I think the setup should be fine

Action: stop the setup

Answer: the final result\\

Begin

Project Structure: \{project\_structure\}

Related Files: \{docs\}
\end{tcolorbox}

\begin{tcolorbox}[colback=gray!5,colframe=black!50,fontupper=\ttfamily,title=Prompt for Verify Agent]
You are a developer. Your task is to verify whether the environment for the given project is set up correctly. Your colleague has set up a Docker environment for the project. You need to verify if it can successfully run the tests of the project.\\

- You interact with a Bash session inside this container.

- The container is based on \{base\_image\}.

- The setup commands that your colleague has run are \{setup\_commands\}

- Project files are located in /testbed within the container, and your current working directory of bash is already set to /testbed.

- Use the same test framework as your colleague, because that aligns with the setup stage.

- Only test commands, skip linting/packaging/publishing commands.

- Do not change the state of the environment, your task is to verify not to fix it. If you see issues, report it not fix it.

- You can tolerate a few test cases failures—as long as most tests pass, it's good enough. \\

\#\# Important Note: \\

Your test command must output detailed pass/fail status for each test item. This is mandatory. For example, with pytest, use the -rA option to get output like: \\

```

PASSED tests/test\_resources.py::test\_fetch\_centromeres

PASSED tests/test\_vis.py::test\_to\_ucsc\_colorstring

```\\

Since we need to parse the test output to extract a test item → status mapping, **this requirement is mandatory**. If you observed that your test command does not produce such detailed output, you must adjust it accordingly. \\

In summary, your goal is:

1. Write the test commands that could output detailed pass/fail status for each test item, you can iterate until it does. (this is mandatory, DO NOT ignore this requirement!!! This is your obligation to correctly identify the test commands to run the test suite of the project, and find a way to output detailed pass/fail status)

2. Run the test command to verify if the environment is set up correctly. If not, report any observed issues. If you think the setup is correct, report none issue.
\end{tcolorbox}

\begin{tcolorbox}[colback=gray!5,colframe=black!50,fontupper=\ttfamily,title=Prompt for Base Image Selection]
Based on related file:

\{related\_files\}

Please recommend a suitable base Docker image. Consider:\\

1. The programming language and version requirements

2. Common system dependencies

3. Use official images when possible \\


Select a base image from the following candidate list:
\{candidate\_images\}

Wrap the image name in a block like <image>python:3.11</image> to indicate your choice.
\end{tcolorbox}



\section{NeurIPS Paper Checklist}

% \begin{itemize}
%     \item You should answer \answerYes{}, \answerNo{}, or \answerNA{}.
%     \item \answerNA{} means either that the question is Not Applicable for that particular paper or the relevant information is Not Available.
%     \item Please provide a short (1–2 sentence) justification right after your answer (even for NA). 
%    % \item {\bf The papers not including the checklist will be desk rejected.}
% \end{itemize}

\begin{enumerate}

\item {\bf Claims}
    \item[] Question: Do the main claims made in the abstract and introduction accurately reflect the paper's contributions and scope?
    \item[] Answer:  \answerYes{} % Replace by \answerYes{}, \answerNo{}, or \answerNA{}.
    \item[] Justification: For detailed contributions and scope, please see Section~\ref{sec:method} and Section~\ref{sec:experiments}%\justificationTODO{}
    \item[] Guidelines:
    \begin{itemize}
        \item The answer NA means that the abstract and introduction do not include the claims made in the paper.
        \item The abstract and/or introduction should clearly state the claims made, including the contributions made in the paper and important assumptions and limitations. A No or NA answer to this question will not be perceived well by the reviewers. 
        \item The claims made should match theoretical and experimental results, and reflect how much the results can be expected to generalize to other settings. 
        \item It is fine to include aspirational goals as motivation as long as it is clear that these goals are not attained by the paper. 
    \end{itemize}

\item {\bf Limitations}
    \item[] Question: Does the paper discuss the limitations of the work performed by the authors?
    \item[] Answer: \answerYes{} % Replace by \answerYes{}, \answerNo{}, or \answerNA{}.
    \item[] Justification: For details about the limitation, please see Section~\ref{sec:limitations} %\justificationTODO{}
    \item[] Guidelines:
    \begin{itemize}
        \item The answer NA means that the paper has no limitation while the answer No means that the paper has limitations, but those are not discussed in the paper. 
        \item The authors are encouraged to create a separate "Limitations" section in their paper.
        \item The paper should point out any strong assumptions and how robust the results are to violations of these assumptions (e.g., independence assumptions, noiseless settings, model well-specification, asymptotic approximations only holding locally). The authors should reflect on how these assumptions might be violated in practice and what the implications would be.
        \item The authors should reflect on the scope of the claims made, e.g., if the approach was only tested on a few datasets or with a few runs. In general, empirical results often depend on implicit assumptions, which should be articulated.
        \item The authors should reflect on the factors that influence the performance of the approach. For example, a facial recognition algorithm may perform poorly when image resolution is low or images are taken in low lighting. Or a speech-to-text system might not be used reliably to provide closed captions for online lectures because it fails to handle technical jargon.
        \item The authors should discuss the computational efficiency of the proposed algorithms and how they scale with dataset size.
        \item If applicable, the authors should discuss possible limitations of their approach to address problems of privacy and fairness.
        \item While the authors might fear that complete honesty about limitations might be used by reviewers as grounds for rejection, a worse outcome might be that reviewers discover limitations that aren't acknowledged in the paper. The authors should use their best judgment and recognize that individual actions in favor of transparency play an important role in developing norms that preserve the integrity of the community. Reviewers will be specifically instructed to not penalize honesty concerning limitations.
    \end{itemize}

\item {\bf Theory assumptions and proofs}
    \item[] Question: For each theoretical result, does the paper provide the full set of assumptions and a complete (and correct) proof?
    \item[] Answer: \answerNA{} % Replace by \answerYes{}, \answerNo{}, or \answerNA{}.
    \item[] Justification: Not applicable as the paper is about dataset and experimental analysis rather than a theory paper. %\justificationTODO{}
    \item[] Guidelines:
    \begin{itemize}
        \item The answer NA means that the paper does not include theoretical results. 
        \item All the theorems, formulas, and proofs in the paper should be numbered and cross-referenced.
        \item All assumptions should be clearly stated or referenced in the statement of any theorems.
        \item The proofs can either appear in the main paper or the supplemental material, but if they appear in the supplemental material, the authors are encouraged to provide a short proof sketch to provide intuition. 
        \item Inversely, any informal proof provided in the core of the paper should be complemented by formal proofs provided in appendix or supplemental material.
        \item Theorems and Lemmas that the proof relies upon should be properly referenced. 
    \end{itemize}

    \item {\bf Experimental result reproducibility}
    \item[] Question: Does the paper fully disclose all the information needed to reproduce the main experimental results of the paper to the extent that it affects the main claims and/or conclusions of the paper (regardless of whether the code and data are provided or not)?
    \item[] Answer: \answerYes{} %\answerTODO{} % Replace by \answerYes{}, \answerNo{}, or \answerNA{}.
    \item[] Justification: We present all the step-wise details in Section~\ref{sec:method} and experimental settings in Section~\ref{sec:experiments}. Besides, we  release all the code, data, and Docker environment for researchers to easily reproduce the results.%\justificationTODO{}
    \item[] Guidelines:
    \begin{itemize}
        \item The answer NA means that the paper does not include experiments.
        \item If the paper includes experiments, a No answer to this question will not be perceived well by the reviewers: Making the paper reproducible is important, regardless of whether the code and data are provided or not.
        \item If the contribution is a dataset and/or model, the authors should describe the steps taken to make their results reproducible or verifiable. 
        \item Depending on the contribution, reproducibility can be accomplished in various ways. For example, if the contribution is a novel architecture, describing the architecture fully might suffice, or if the contribution is a specific model and empirical evaluation, it may be necessary to either make it possible for others to replicate the model with the same dataset, or provide access to the model. In general. releasing code and data is often one good way to accomplish this, but reproducibility can also be provided via detailed instructions for how to replicate the results, access to a hosted model (e.g., in the case of a large language model), releasing of a model checkpoint, or other means that are appropriate to the research performed.
        \item While NeurIPS does not require releasing code, the conference does require all submissions to provide some reasonable avenue for reproducibility, which may depend on the nature of the contribution. For example
        \begin{enumerate}
            \item If the contribution is primarily a new algorithm, the paper should make it clear how to reproduce that algorithm.
            \item If the contribution is primarily a new model architecture, the paper should describe the architecture clearly and fully.
            \item If the contribution is a new model (e.g., a large language model), then there should either be a way to access this model for reproducing the results or a way to reproduce the model (e.g., with an open-source dataset or instructions for how to construct the dataset).
            \item We recognize that reproducibility may be tricky in some cases, in which case authors are welcome to describe the particular way they provide for reproducibility. In the case of closed-source models, it may be that access to the model is limited in some way (e.g., to registered users), but it should be possible for other researchers to have some path to reproducing or verifying the results.
        \end{enumerate}
    \end{itemize}


\item {\bf Open access to data and code}
    \item[] Question: Does the paper provide open access to the data and code, with sufficient instructions to faithfully reproduce the main experimental results, as described in supplemental material?
    \item[] Answer: \answerYes{} %\answerTODO{} % Replace by \answerYes{}, \answerNo{}, or \answerNA{}.
    \item[] Justification: All the code and data are open accessible with proper README instructions.
    \item[] Guidelines:
    \begin{itemize}
        \item The answer NA means that paper does not include experiments requiring code.
        \item Please see the NeurIPS code and data submission guidelines (\url{https://nips.cc/public/guides/CodeSubmissionPolicy}) for more details.
        \item While we encourage the release of code and data, we understand that this might not be possible, so “No” is an acceptable answer. Papers cannot be rejected simply for not including code, unless this is central to the contribution (e.g., for a new open-source benchmark).
        \item The instructions should contain the exact command and environment needed to run to reproduce the results. See the NeurIPS code and data submission guidelines (\url{https://nips.cc/public/guides/CodeSubmissionPolicy}) for more details.
        \item The authors should provide instructions on data access and preparation, including how to access the raw data, preprocessed data, intermediate data, and generated data, etc.
        \item The authors should provide scripts to reproduce all experimental results for the new proposed method and baselines. If only a subset of experiments are reproducible, they should state which ones are omitted from the script and why.
        \item At submission time, to preserve anonymity, the authors should release anonymized versions (if applicable).
        \item Providing as much information as possible in supplemental material (appended to the paper) is recommended, but including URLs to data and code is permitted.
    \end{itemize}


\item {\bf Experimental setting/details}
    \item[] Question: Does the paper specify all the training and test details (e.g., data splits, hyperparameters, how they were chosen, type of optimizer, etc.) necessary to understand the results?
    \item[] Answer:\answerYes{} % \answerTODO{} % Replace by \answerYes{}, \answerNo{}, or \answerNA{}.
    \item[] Justification: We presented the experimental details in Appendix~\ref{sec:experimental_setup} 
    \item[] Guidelines:
    \begin{itemize}
        \item The answer NA means that the paper does not include experiments.
        \item The experimental setting should be presented in the core of the paper to a level of detail that is necessary to appreciate the results and make sense of them.
        \item The full details can be provided either with the code, in appendix, or as supplemental material.
    \end{itemize}

\item {\bf Experiment statistical significance}
    \item[] Question: Does the paper report error bars suitably and correctly defined or other appropriate information about the statistical significance of the experiments?
    \item[] Answer: \answerNo{} % Replace by \answerYes{}, \answerNo{}, or \answerNA{}.
    \item[] Justification: We do not report the experiment statistical significance due to the limited budget for repeat experiments. We have controlled the experimental randomness by setting temperature to 0.
    \item[] Guidelines:
    \begin{itemize}
        \item The answer NA means that the paper does not include experiments.
        \item The authors should answer "Yes" if the results are accompanied by error bars, confidence intervals, or statistical significance tests, at least for the experiments that support the main claims of the paper.
        \item The factors of variability that the error bars are capturing should be clearly stated (for example, train/test split, initialization, random drawing of some parameter, or overall run with given experimental conditions).
        \item The method for calculating the error bars should be explained (closed form formula, call to a library function, bootstrap, etc.)
        \item The assumptions made should be given (e.g., Normally distributed errors).
        \item It should be clear whether the error bar is the standard deviation or the standard error of the mean.
        \item It is OK to report 1-sigma error bars, but one should state it. The authors should preferably report a 2-sigma error bar than state that they have a 96\% CI, if the hypothesis of Normality of errors is not verified.
        \item For asymmetric distributions, the authors should be careful not to show in tables or figures symmetric error bars that would yield results that are out of range (e.g. negative error rates).
        \item If error bars are reported in tables or plots, The authors should explain in the text how they were calculated and reference the corresponding figures or tables in the text.
    \end{itemize}

\item {\bf Experiments compute resources}
    \item[] Question: For each experiment, does the paper provide sufficient information on the computer resources (type of compute workers, memory, time of execution) needed to reproduce the experiments?
    \item[] Answer: \answerYes{} % Replace by \answerYes{}, \answerNo{}, or \answerNA{}.
    \item[] Justification: We provide details about the experimental settings in Section~\ref{sec:experiments} and Appendix~\ref{sec:experimental_setup}.
    \item[] Guidelines:
    \begin{itemize}
        \item The answer NA means that the paper does not include experiments.
        \item The paper should indicate the type of compute workers CPU or GPU, internal cluster, or cloud provider, including relevant memory and storage.
        \item The paper should provide the amount of compute required for each of the individual experimental runs as well as estimate the total compute. 
        \item The paper should disclose whether the full research project required more compute than the experiments reported in the paper (e.g., preliminary or failed experiments that didn't make it into the paper). 
    \end{itemize}
    
\item {\bf Code of ethics}
    \item[] Question: Does the research conducted in the paper conform, in every respect, with the NeurIPS Code of Ethics \url{https://neurips.cc/public/EthicsGuidelines}?
    \item[] Answer: \answerYes{} % Replace by \answerYes{}, \answerNo{}, or \answerNA{}.
    \item[] Justification: We have followed all the necessary code of ethics to conduct the research % \justificationTODO{}
    \item[] Guidelines:
    \begin{itemize}
        \item The answer NA means that the authors have not reviewed the NeurIPS Code of Ethics.
        \item If the authors answer No, they should explain the special circumstances that require a deviation from the Code of Ethics.
        \item The authors should make sure to preserve anonymity (e.g., if there is a special consideration due to laws or regulations in their jurisdiction).
    \end{itemize}


\item {\bf Broader impacts}
    \item[] Question: Does the paper discuss both potential positive societal impacts and negative societal impacts of the work performed?
    \item[] Answer: \answerNA{} % Replace by \answerYes{}, \answerNo{}, or \answerNA{}.
    \item[] Justification: The paper mainly focuses on issue fixing, a specific software engineering task, which has no relation with societal issues.%\justificationTODO{}
    \item[] Guidelines:
    \begin{itemize}
        \item The answer NA means that there is no societal impact of the work performed.
        \item If the authors answer NA or No, they should explain why their work has no societal impact or why the paper does not address societal impact.
        \item Examples of negative societal impacts include potential malicious or unintended uses (e.g., disinformation, generating fake profiles, surveillance), fairness considerations (e.g., deployment of technologies that could make decisions that unfairly impact specific groups), privacy considerations, and security considerations.
        \item The conference expects that many papers will be foundational research and not tied to particular applications, let alone deployments. However, if there is a direct path to any negative applications, the authors should point it out. For example, it is legitimate to point out that an improvement in the quality of generative models could be used to generate deepfakes for disinformation. On the other hand, it is not needed to point out that a generic algorithm for optimizing neural networks could enable people to train models that generate Deepfakes faster.
        \item The authors should consider possible harms that could arise when the technology is being used as intended and functioning correctly, harms that could arise when the technology is being used as intended but gives incorrect results, and harms following from (intentional or unintentional) misuse of the technology.
        \item If there are negative societal impacts, the authors could also discuss possible mitigation strategies (e.g., gated release of models, providing defenses in addition to attacks, mechanisms for monitoring misuse, mechanisms to monitor how a system learns from feedback over time, improving the efficiency and accessibility of ML).
    \end{itemize}
    
\item {\bf Safeguards}
    \item[] Question: Does the paper describe safeguards that have been put in place for responsible release of data or models that have a high risk for misuse (e.g., pretrained language models, image generators, or scraped datasets)?
    \item[] Answer: \answerNA{} %\answerTODO{} % Replace by \answerYes{}, \answerNo{}, or \answerNA{}.
    \item[] Justification: We do not release models. The benchmark are about real-world issue resolving and we provided Docker containers as the isolated environment, thereby eliminating the safety risks.
    \item[] Guidelines:
    \begin{itemize}
        \item The answer NA means that the paper poses no such risks.
        \item Released models that have a high risk for misuse or dual-use should be released with necessary safeguards to allow for controlled use of the model, for example by requiring that users adhere to usage guidelines or restrictions to access the model or implementing safety filters. 
        \item Datasets that have been scraped from the Internet could pose safety risks. The authors should describe how they avoided releasing unsafe images.
        \item We recognize that providing effective safeguards is challenging, and many papers do not require this, but we encourage authors to take this into account and make a best faith effort.
    \end{itemize}

\item {\bf Licenses for existing assets}
    \item[] Question: Are the creators or original owners of assets (e.g., code, data, models), used in the paper, properly credited and are the license and terms of use explicitly mentioned and properly respected?
    \item[] Answer: \answerYes{} %\answerTODO{} % Replace by \answerYes{}, \answerNo{}, or \answerNA{}.
    \item[] Justification: All the code, data or models used in this work are properly cited as References.
    \item[] Guidelines:
    \begin{itemize}
        \item The answer NA means that the paper does not use existing assets.
        \item The authors should cite the original paper that produced the code package or dataset.
        \item The authors should state which version of the asset is used and, if possible, include a URL.
        \item The name of the license (e.g., CC-BY 4.0) should be included for each asset.
        \item For scraped data from a particular source (e.g., website), the copyright and terms of service of that source should be provided.
        \item If assets are released, the license, copyright information, and terms of use in the package should be provided. For popular datasets, \url{paperswithcode.com/datasets} has curated licenses for some datasets. Their licensing guide can help determine the license of a dataset.
        \item For existing datasets that are re-packaged, both the original license and the license of the derived asset (if it has changed) should be provided.
        \item If this information is not available online, the authors are encouraged to reach out to the asset's creators.
    \end{itemize}

\item {\bf New assets}
    \item[] Question: Are new assets introduced in the paper well documented and is the documentation provided alongside the assets?
    \item[] Answer: \answerYes{}% Replace by \answerYes{}, \answerNo{}, or \answerNA{}.
    \item[] Justification: The benchmark, datasets and code are well-documented.
    \item[] Guidelines:
    \begin{itemize}
        \item The answer NA means that the paper does not release new assets.
        \item Researchers should communicate the details of the dataset/code/model as part of their submissions via structured templates. This includes details about training, license, limitations, etc. 
        \item The paper should discuss whether and how consent was obtained from people whose asset is used.
        \item At submission time, remember to anonymize your assets (if applicable). You can either create an anonymized URL or include an anonymized zip file.
    \end{itemize}

\item {\bf Crowdsourcing and research with human subjects}
    \item[] Question: For crowdsourcing experiments and research with human subjects, does the paper include the full text of instructions given to participants and screenshots, if applicable, as well as details about compensation (if any)? 
    \item[] Answer: \answerNA{} % Replace by \answerYes{}, \answerNo{}, or \answerNA{}.
    \item[] Justification: the paper does not involve crowdsourcing nor research with human subjects.
    \item[] Guidelines:
    \begin{itemize}
        \item The answer NA means that the paper does not involve crowdsourcing nor research with human subjects.
        \item Including this information in the supplemental material is fine, but if the main contribution of the paper involves human subjects, then as much detail as possible should be included in the main paper. 
        \item According to the NeurIPS Code of Ethics, workers involved in data collection, curation, or other labor should be paid at least the minimum wage in the country of the data collector. 
    \end{itemize}

\item {\bf Institutional review board (IRB) approvals or equivalent for research with human subjects}
    \item[] Question: Does the paper describe potential risks incurred by study participants, whether such risks were disclosed to the subjects, and whether Institutional Review Board (IRB) approvals (or an equivalent approval/review based on the requirements of your country or institution) were obtained?
    \item[] Answer: \answerNA{}% Replace by \answerYes{}, \answerNo{}, or \answerNA{}.
    \item[] Justification: the paper does not involve crowdsourcing nor research with human subjects
    \item[] Guidelines:
    \begin{itemize}
        \item The answer NA means that the paper does not involve crowdsourcing nor research with human subjects.
        \item Depending on the country in which research is conducted, IRB approval (or equivalent) may be required for any human subjects research. If you obtained IRB approval, you should clearly state this in the paper. 
        \item We recognize that the procedures for this may vary significantly between institutions and locations, and we expect authors to adhere to the NeurIPS Code of Ethics and the guidelines for their institution. 
        \item For initial submissions, do not include any information that would break anonymity (if applicable), such as the institution conducting the review.
    \end{itemize}

\item {\bf Declaration of LLM usage}
    \item[] Question: Does the paper describe the usage of LLMs if it is an important, original, or non-standard component of the core methods in this research? Note that if the LLM is used only for writing, editing, or formatting purposes and does not impact the core methodology, scientific rigorousness, or originality of the research, declaration is not required.
    %this research? 
    \item[] Answer: \answerYes{} % Replace by \answerYes{}, \answerNo{}, or \answerNA{}.
    \item[] Justification: LLMs are core part in the paper, we have listed all the LLM models used in our experiments in Section~\ref{sec:experiments}.
    \item[] Guidelines:
    \begin{itemize}
        \item The answer NA means that the core method development in this research does not involve LLMs as any important, original, or non-standard components.
        \item Please refer to our LLM policy (\url{https://neurips.cc/Conferences/2025/LLM}) for what should or should not be described.
    \end{itemize}

\end{enumerate}
